\chapter{Hệ thống đề xuất}
\label{chap:HeThongDeXuat}

\section{Mô tả hệ thống}
\label{sec:MoTaHeThong}

Hệ thống được đề xuất là một \textbf{hệ thống quản lý thư viện trực tuyến tích hợp AI hỗ trợ tìm kiếm và gợi ý sách}. Mục tiêu cốt lõi của dự án là xây dựng một nền tảng hiện đại, kết hợp hài hòa giữa hai yếu tố: (1) cung cấp đầy đủ các chức năng để quản lý thư viện một cách hiệu quả, và (2) ứng dụng Trí tuệ nhân tạo (AI) để mang lại trải nghiệm thông minh cho người dùng cuối. Thay vì chỉ dừng lại ở mức tra cứu từ khóa cơ bản như các hệ thống hiện tại, hệ thống đề xuất hướng tới việc trở thành một trợ lý ảo, giúp người dùng khám phá và tiếp cận nguồn tài nguyên tri thức một cách tối ưu nhất.

Hệ thống được thiết kế để phục vụ hai nhóm người dùng chính, mỗi nhóm có vai trò và nhu- cầu riêng biệt:
\begin{itemize}
    \item \textbf{Bạn đọc (Sinh viên, Giảng viên, Nghiên cứu viên):} Là nhóm người dùng cuối, trực tiếp khai thác tài nguyên của thư viện. Nhu cầu chính của họ là tìm kiếm tài liệu nhanh chóng, đọc tóm tắt để đánh giá nội dung, thực hiện mượn sách và nhận được các gợi ý tài liệu liên quan để phục vụ cho việc học tập và nghiên cứu.
    
    \item \textbf{Thủ thư (Người quản lý thư viện):} Là nhóm người dùng quản trị, chịu trách nhiệm vận hành thư viện. Vai trò của họ bao gồm việc quản lý kho tài liệu, quản lý thông tin bạn đọc, theo dõi các hoạt động mượn - trả, và sử dụng các công cụ thống kê, báo cáo tự động để đánh giá hiệu quả hoạt động của thư viện.
\end{itemize}

\section{Hướng tiếp cận thiết kế}
\label{subsec:DesignApproach}

Hệ thống Quản lý Thư viện sẽ được thiết kế theo phương pháp ``từ dưới lên'' (bottom-up design). Hướng tiếp cận này tập trung vào việc xây dựng và kiểm thử các thành phần cơ bản, độc lập trước khi tích hợp chúng lại thành một hệ thống lớn hơn. Quá trình này bao gồm các bước sau:
\begin{enumerate}
    \item Xác định và thiết kế các thành phần nhỏ, cốt lõi nhất của hệ thống (ví dụ: các module quản lý sách, quản lý người dùng).
    \item Xây dựng và kiểm thử riêng lẻ từng thành phần để đảm bảo chúng hoạt động chính xác.
    \item Tích hợp các thành phần nhỏ đã hoàn thiện để xây dựng các chức năng hoặc hệ thống con lớn hơn.
    \item Lặp lại quy trình này cho đến khi hoàn thành thiết kế của toàn bộ hệ thống.
\end{enumerate}
\newpage
\section{Yêu cầu}
\label{sec:YeuCau}

\subsection{Yêu cầu chức năng}
\label{subsec:YeuCauChucNang}
Các yêu cầu chức năng của hệ thống được đặc tả chi tiết trong bảng dưới đây, phân loại theo từng nhóm chức năng và người dùng tương tác.

\begin{longtable}{|p{0.1\textwidth}|p{0.22\textwidth}|p{0.13\textwidth}|p{0.45\textwidth}|}
\caption{Bảng đặc tả yêu cầu hệ thống} \label{tab:YeuCauHeThong} \\
\hline
\textbf{ID} & \textbf{Tên yêu cầu} & \textbf{Loại} & \textbf{Mô tả chi tiết} \\
\hline
\endfirsthead
\multicolumn{4}{c}%
{{\bfseries \tablename\ \thetable{} -- tiếp theo}} \\
\hline
\textbf{ID} & \textbf{Tên yêu cầu} & \textbf{Loại} & \textbf{Mô tả chi tiết} \\
\hline
\endhead
\hline
\endfoot
%-- Yêu cầu về Dữ liệu --%
\textbf{R1} & Quản lý thực thể dữ liệu & Dữ liệu & Hệ thống phải có khả năng lưu trữ và quản lý ba thực thể dữ liệu cốt lõi: \textbf{Sách}, \textbf{Người dùng}, và \textbf{Lịch sử Giao dịch} (mượn/trả). \\
\hline
\textbf{R2} & Phân biệt Đầu sách và Bản sao & Dữ liệu & Hệ thống phải phân biệt giữa ``Đầu sách'' và ``Bản sao''. Mỗi bản sao phải có một mã định danh duy nhất (barcode), thông tin trạng thái và \textbf{vị trí kệ sách} để xác định vị trí vật lý. \\
\hline
\textbf{R3} & Thuộc tính Đầu sách & Dữ liệu & Mỗi đầu sách phải có các thuộc tính cơ bản bao gồm: \textbf{ISBN}, tiêu đề, tác giả, thể loại và ngày xuất bản. \\
\hline
%-- Yêu cầu về Người dùng và Nghiệp vụ --%
\textbf{R4} & Quản lý vai trò người dùng & Chức năng & Hệ thống phải hỗ trợ hai vai trò người dùng: \textbf{Thủ thư} (quản trị) và \textbf{Bạn đọc} (thành viên). Mỗi người dùng phải được định danh bởi một mã thẻ thư viện duy nhất. \\
\hline
\textbf{R5} & Cấu hình chính sách mượn/trả & Nghiệp vụ & Hệ thống phải cho phép Thủ thư cấu hình các \textbf{chính sách mượn trả}, ví dụ: giới hạn số lượng sách mượn tối đa là \textbf{5 cuốn} và thời gian mượn tiêu chuẩn là \textbf{14 ngày}. \\
\hline
%-- Yêu cầu về Chức năng --%
\textbf{R6} & Tìm kiếm ngữ nghĩa & Chức năng & Hệ thống phải cung cấp chức năng \textbf{tìm kiếm ngữ nghĩa}, cho phép tìm kiếm sách theo tiêu đề, tác giả, hoặc chủ đề, và xử lý được các truy vấn tiếng Việt. \\
\hline
\textbf{R7} & Chức năng Đặt chỗ & Chức năng & Bạn đọc phải có khả năng \textbf{đặt chỗ} một đầu sách khi không có bản sao nào sẵn có. Mỗi đầu sách có thể được đặt chỗ bởi \textbf{nhiều thành viên} tại một thời điểm. \\
\hline
\textbf{R8} & Chức năng Gia hạn & Chức năng & Hệ thống phải cho phép Bạn đọc \textbf{gia hạn} thời gian mượn sách, với điều kiện sách đó không đang được người khác đặt chỗ. \\
\hline
\textbf{R9} & Hệ thống thông báo tự động & Chức năng & Hệ thống phải \textbf{tự động gửi thông báo} cho Bạn đọc về các sự kiện quan trọng, bao gồm: sách sắp đến hạn trả, sách đã quá hạn, và sách đặt chỗ đã có sẵn. \\
\hline
\end{longtable}

\subsection{Yêu cầu phi chức năng}
\label{subsec:YeuCauPhiChucNang}
Hệ thống cần đảm bảo các yêu cầu phi chức năng trọng yếu sau:
\begin{itemize}
    \item \textbf{Tính khả dụng (Usability):} Giao diện người dùng phải nhất quán, trực quan và dễ học cho cả bạn đọc và thủ thư. Hệ thống phải có thiết kế đáp ứng (responsive), hoạt động tốt trên các trình duyệt web phổ biến và các kích thước màn hình khác nhau.
    
    \item \textbf{Hiệu năng (Performance):} Các chức năng tương tác phải có thời gian phản hồi nhanh. Cụ thể, thời gian tải trang chính không quá 3 giây và thời gian trả về kết quả của chức năng tìm kiếm và gợi ý không quá 2 giây trong điều kiện tải thông thường.
    
    \item \textbf{Độ tin cậy (Reliability):} Hệ thống phải hoạt động ổn định, sẵn sàng 24/7 với thời gian hoạt động (uptime) đạt tối thiểu 99.5\%. Dữ liệu phải được sao lưu định kỳ để đảm bảo khả năng phục hồi khi có sự cố.
    
    \item \textbf{Bảo mật (Security):} Dữ liệu cá nhân của người dùng phải được mã hóa. Hệ thống cần có cơ chế xác thực và phân quyền rõ ràng, đảm bảo bạn đọc không thể truy cập các chức năng của thủ thư và ngược lại.
    
    \item \textbf{Khả năng bảo trì và mở rộng (Maintainability \& Scalability):} Mã nguồn phải được viết rõ ràng, có tài liệu và tuân theo các chuẩn thiết kế. Kiến trúc tách biệt backend-frontend cho phép các nhóm phát triển độc lập và dễ dàng nâng cấp, tích hợp các tính năng mới trong tương lai.
\end{itemize}

\subsection{Yêu cầu dữ liệu}
\label{subsec:YeuCauDuLieu}
Hệ thống yêu cầu và quản lý các nhóm dữ liệu chính sau:
\begin{itemize}
    \item \textbf{Dữ liệu lõi (Core Data):} Bao gồm thông tin chi tiết về sách (metadata), tác giả, thể loại và thông tin tài khoản của người dùng (bạn đọc, thủ thư).
    
    \item \textbf{Dữ liệu nghiệp vụ (Transactional Data):} Gồm các bản ghi về hoạt động mượn và trả sách, lịch sử giao dịch, và thông tin về các khoản phí phạt (nếu có).
    
    \item \textbf{Dữ liệu hành vi (Behavioral Data):} Gồm lịch sử các truy vấn tìm kiếm, các lượt xem chi tiết sách, và các tương tác khác của người dùng trên hệ thống. Đây là nguồn dữ liệu quan trọng cho các thuật toán gợi ý.
\end{itemize}

\section{Kết chương}
\label{sec:KetChuong4}

Chương này đã trình bày một cách tổng quan và chi tiết về hệ thống quản lý thư viện thông minh được đề xuất. Hệ thống đã được mô tả rõ ràng về mục tiêu, định hướng và hai nhóm người dùng chính là Bạn đọc và Thủ thư.

Phần trọng tâm của chương đã tập trung đặc tả các yêu cầu cốt lõi của hệ thống. Các \textbf{yêu cầu chức năng} đã được liệt kê chi tiết, bao quát từ các nghiệp vụ quản lý cơ bản cho đến các tính năng AI nâng cao như tìm kiếm ngữ nghĩa và gợi ý cá nhân hóa. Đồng thời, các \textbf{yêu cầu phi chức năng} quan trọng như hiệu năng, tính khả dụng, và bảo mật cũng đã được xác định để đảm bảo hệ thống hoạt động hiệu quả và an toàn. Cuối cùng, các \textbf{yêu cầu về dữ liệu} cần thiết cho việc vận hành và huấn luyện các mô hình AI cũng đã được làm rõ.

Những yêu cầu này tạo thành một nền tảng vững chắc, là cơ sở để tiến hành quá trình phân tích và thiết kế chi tiết hệ thống sẽ được trình bày ở Chương 5.